% practice_answers.tex -- Module 5 (Advanced Mechanics): worked answers to the Practice set.
% One practiceans per practiceq in practice.tex, in the same order (24 answers).
% Second, independent recalculation of every final answer:
% Third check (every question re-solved before reading these answers):
% MC and calculation answers keep their first lines together with \Needspace / \nobreak
% (conditional, layout-independent; no \clearpage or \pagebreak).
\begin{practiceanswers}

% ===================================================================== Section I
% ---- Q1
\Needspace{9\baselineskip}
\begin{practiceans}
  \correct{D}
  At the top, the vertical velocity is zero, but the horizontal velocity $u_x$ is still
  there. So the velocity is horizontal. Gravity still acts, so the acceleration is
  \qty{9.8}{m.s^{-2}} downwards at every point of the path, including the top.
  \begin{itemize}
    \item \xmark\ A: the velocity is not zero ($v_x$ is still there), and gravity does not
          switch off at the top.
    \item \xmark\ B: only the \emph{vertical} part of the velocity is zero.
    \item \xmark\ C: gravity still acts at the top, so the acceleration is not zero.
  \end{itemize}
\end{practiceans}

% ---- Q2
\Needspace{9\baselineskip}
\begin{practiceans}
  \correct{A}
  $v_y = u_y + a_y t$ with $a_y = \qty{-9.8}{m.s^{-2}}$. This is a straight line with a
  gradient of $-9.8$. It starts at $+u_y$, crosses zero at the top of the flight (halfway),
  and ends at $-u_y$ when the ball lands (same speed, now downwards).
  \begin{itemize}
    \item \xmark\ B: $v_y$ never becomes negative here, so the ball would never come down. This is
          the graph of the \emph{size} of $v_y$, not of the velocity.
    \item \xmark\ C: a constant value is the graph of the horizontal velocity $v_x$.
    \item \xmark\ D: this curve is the shape of the height--time graph, not the velocity--time
          graph.
  \end{itemize}
\end{practiceans}

% ---- Q3
\Needspace{9\baselineskip}
\begin{practiceans}
  \correct{C}
  $F_\mathrm{c} = \dfrac{mv^2}{r}$. Twice the speed makes $v^2$ bigger by $2^2 = 4$. Half
  the radius doubles the force, because we divide by a number half as big. Together:
  $4 \times 2 = 8$, so the force is $8F$.
  \begin{itemize}
    \item \xmark\ A ($2F$): multiplies by the new radius instead of dividing:
          $4 \times \tfrac{1}{2} = 2$.
    \item \xmark\ B ($4F$): forgets to square the speed: $2 \times 2 = 4$.
    \item \xmark\ D ($16F$): squares the change in radius as well: $4 \times 2^2 = 16$.
  \end{itemize}
\end{practiceans}

% ---- Q4
\Needspace{9\baselineskip}
\begin{practiceans}
  \correct{D}
  The normal force acts at \ang{90} to the road. On a banked road it tilts towards the centre
  of the curve. Its vertical component balances the weight. Its horizontal component points
  to the centre, so it provides the centripetal force.
  \begin{itemize}
    \item \xmark\ A: the weight acts straight down, but the centre of the circle is
          sideways (horizontal) from the car.
    \item \xmark\ B: at the design speed no sideways friction is needed. Friction is needed
          only at other speeds.
    \item \xmark\ C: the vertical component balances the weight. It has no part towards the
          centre.
  \end{itemize}
\end{practiceans}

% ---- Q5
\Needspace{9\baselineskip}
\begin{practiceans}
  \correct{C}
  Every point on the merry-go-round turns through the same angle in the same time, so P and Q
  have the same $\omega$. With $v = \omega r$, Q (twice the radius) moves twice as fast.
  Then $a_\mathrm{c} = \dfrac{v^2}{r} = \dfrac{(\omega r)^2}{r} = \omega^2 r$. With the same
  $\omega$, $a_\mathrm{c}$ is proportional to $r$, so Q has twice the acceleration of P.
  \begin{itemize}
    \item \xmark\ A: $v = \omega r$, so Q moves twice as fast as P.
    \item \xmark\ B: $\omega$ is the same for every point on the merry-go-round.
    \item \xmark\ D: Q's $v^2$ is 4 times bigger, but its $r$ is also 2 times bigger:
          $4 \div 2 = 2$, not 4.
  \end{itemize}
\end{practiceans}

% ---- Q6
\Needspace{9\baselineskip}
\begin{practiceans}
  \correct{B}
  Use $\tau = rF\sin\theta$ for each push ($\theta$ = angle between the door and the force):
  \begin{center}
  \begin{tabular}{@{}cl@{}}
    A & $0.90 \times 30 \times \sin\ang{90} = \qty{27}{N.m}$\\
    B & $0.60 \times 50 \times \sin\ang{90} = \qty{30}{N.m}$ \quad (largest)\\
    C & $0.90 \times 80 \times \sin\ang{20} = \qty{24.6}{N.m}$\\
    D & $0.25 \times 100 \times \sin\ang{90} = \qty{25}{N.m}$
  \end{tabular}
  \end{center}
  \begin{itemize}
    \item \xmark\ A: the push is far from the hinge, but the force is small.
    \item \xmark\ C: the force is large and far from the hinge, but at \ang{20} to the door
          only the part $F\sin\ang{20}$ turns it. (Using $\cos\ang{20}$ by mistake gives
          \qty{67.7}{N.m} and makes C look largest.)
    \item \xmark\ D: the force is the largest, but it acts very close to the hinge.
  \end{itemize}
\end{practiceans}

% ---- Q7
\Needspace{9\baselineskip}
\begin{practiceans}
  \correct{B}
  $g = \dfrac{GM}{r^2}$. For the planet,
  $g_\mathrm{planet} = \dfrac{G(4M_\mathrm{E})}{(2r_\mathrm{E})^2}
   = \dfrac{4}{4} \times \dfrac{GM_\mathrm{E}}{r_\mathrm{E}^2} = 1 \times 9.8
   = \qty{9.8}{N.kg^{-1}}$. The bigger mass and the bigger radius cancel out.
  \begin{itemize}
    \item \xmark\ A (2.45): uses the bigger radius but forgets the bigger mass:
          $9.8 \div 4$.
    \item \xmark\ C (19.6): forgets to square the radius: $9.8 \times 4 \div 2$.
    \item \xmark\ D (39.2): uses the bigger mass but forgets the bigger radius: $9.8 \times 4$.
  \end{itemize}
\end{practiceans}

% ---- Q8
\Needspace{9\baselineskip}
\begin{practiceans}
  \correct{A}
  In a circular orbit, $K = \dfrac{GMm}{2r}$, $U = -\dfrac{GMm}{r}$ and
  $E = -\dfrac{GMm}{2r}$. When $r$ gets bigger, $K$ gets smaller (the satellite moves more
  slowly). $U$ and $E$ become less negative, so they increase. The engine supplies the extra
  energy.
  \begin{itemize}
    \item \xmark\ B: a higher orbit has more total energy. It must come from the engine, so $E$
          cannot stay the same.
    \item \xmark\ C: these are the changes for a move to a \emph{lower} orbit.
    \item \xmark\ D: $v = \sqrt{\dfrac{GM}{r}}$ is smaller in a higher orbit, so $K$ decreases.
  \end{itemize}
\end{practiceans}

% ===================================================================== Section II
% ---- Q9 drone
\begin{practiceans}
  \given{$u_x = \qty{7.5}{m.s^{-1}}$; $u_y = 0$ (the kit starts with the drone's horizontal
    velocity); $s_y = \qty{-20}{m}$ (up is positive); $a_y = \qty{-9.8}{m.s^{-2}}$}
  \find{(a) the time $t$ \quad (b) the horizontal distance $s_x$ \quad (c) where the drone is}
  \formula{$s_y = u_y t + \frac{1}{2}a_y t^2$ \quad and \quad $s_x = u_x t$}
  \substitute{\textbf{(a)} Vertical motion:\par\nobreak
    \eqlines{\qty{-20}{m} &= 0 + \tfrac{1}{2}(\qty{-9.8}{m.s^{-2}})\,t^2 \quad\text{($u_y = 0$)}\\
      t^2 &= \frac{\qty{-20}{m}}{\qty{-4.9}{m.s^{-2}}} = \qty{4.082}{s^2} \quad\text{(divide both sides by $-4.9$)}\\
      t &= \sqrt{\qty{4.082}{s^2}} = \qty{2.020}{s} \quad\text{(square root)}}\par
    \textbf{(b)} Horizontal motion, same $t$:
    $s_x = \qty{7.5}{m.s^{-1}} \times \qty{2.020}{s} = \qty{15.15}{m}$.\par
    \textbf{(c)} No horizontal force acts on the kit. So its horizontal velocity stays
    \qty{7.5}{m.s^{-1}}, the same as the drone's. In \qty{2.0}{s} both move \qty{15}{m}
    forwards.}
  \Needspace{5\baselineskip}% layout QA: keep Answer and Check together
  \answer{(a) \final{$t = \qty{2.0}{s}$} \quad (b) \final{$s_x = \qty{15}{m}$} \quad
    (c) the drone is \textbf{directly above} the kit when it lands.}
  \checkstep{$\tfrac{1}{2} \times 9.8 \times 2.020^2 = \qty{20.0}{m}$, the drop height \cmark.}
  \begin{markingcriteria}
    \criterion{(a) Correct time, \qty{2.0}{s}, with working}{1}
    \criterion{(b) Correct horizontal distance, \qty{15}{m}}{1}
    \criterion{(c) Directly above the kit, because both keep the same horizontal velocity}{1}
  \end{markingcriteria}
\end{practiceans}

% ---- Q10 centrifuge
\begin{practiceans}
  \given{$r = \qty{8.0}{m}$, \ $m = \qty{72}{kg}$, \
    $a_\mathrm{c} = 6.0 \times 9.8 = \qty{58.8}{m.s^{-2}}$}
  \find{(a) the speed $v$ \quad (b) the period $T$ \quad (c) the centripetal force
    $F_\mathrm{c}$}
  \formula{$a_\mathrm{c} = \dfrac{v^2}{r}$, \quad $v = \dfrac{2\pi r}{T}$, \quad
    $F_\mathrm{c} = \dfrac{mv^2}{r} = ma_\mathrm{c}$}
  \substitute{\eqlines{v^2 &= a_\mathrm{c}\, r = \qty{58.8}{m.s^{-2}} \times \qty{8.0}{m}
        = \qty{470.4}{m^2.s^{-2}} \quad\text{(multiply both sides by $r$)}\\
      v &= \sqrt{\qty{470.4}{m^2.s^{-2}}} = \qty{21.69}{m.s^{-1}} \quad\text{(square root)}\\
      T &= \frac{2\pi r}{v} = \frac{2\pi \times \qty{8.0}{m}}{\qty{21.69}{m.s^{-1}}}
        = \frac{\qty{50.27}{m}}{\qty{21.69}{m.s^{-1}}} = \qty{2.318}{s}\\
      F_\mathrm{c} &= ma_\mathrm{c} = \qty{72}{kg} \times \qty{58.8}{m.s^{-2}}
        = \qty{4234}{N}}}
  \answer{(a) \final{$v = \qty{22}{m.s^{-1}}$} \quad (b) \final{$T = \qty{2.3}{s}$} \quad
    (c) \final{$F_\mathrm{c} = \qty{4.2e3}{N}$}, towards the centre}
  \checkstep{the pilot's weight is $72 \times 9.8 = \qty{706}{N}$, and
    $4234 \div 706 = 6.0$, as the question says \cmark.}
  \begin{markingcriteria}
    \criterion{(a) Correct speed, \qty{22}{m.s^{-1}}}{1}
    \criterion{(b) Correct period, \qty{2.3}{s}}{1}
    \criterion{(c) Correct force, \qty{4.2e3}{N}}{1}
  \end{markingcriteria}
\end{practiceans}

% ---- Q11 wind turbine
\begin{practiceans}
  \given{$r = \qty{45.0}{m}$ (tip); 16.0 turns in \qty{60}{s}; one turn $= 2\pi$ rad}
  \find{(a) $\omega$ \quad (b) the tip speed $v$ in \si{km.h^{-1}} \quad (c) $\omega$ and $v$
    at M}
  \formula{$\omega = \dfrac{\Delta\theta}{t}$ \quad and \quad $v = \omega r$ (not on the
    formula sheet; it comes from $v = \dfrac{2\pi r}{T}$)}
  \substitute{\eqlines{\Delta\theta &= 16.0 \times 2\pi = \qty{100.5}{rad} \quad\text{(turns to radians)}\\
      \omega &= \frac{\qty{100.5}{rad}}{\qty{60}{s}} = \qty{1.676}{rad.s^{-1}}\\
      v &= \omega r = \qty{1.676}{rad.s^{-1}} \times \qty{45.0}{m} = \qty{75.40}{m.s^{-1}}\\
      v &= 75.40 \times 3.6 = \qty{271.4}{km.h^{-1}} \quad\text{(\si{m.s^{-1}} to \si{km.h^{-1}}: $\times 3.6$)}}\par
    \textbf{(c)} M turns through the same angle in the same time as the tip, so it has the
    same $\omega$. Its radius is \qty{22.5}{m}, so
    $v = \qty{1.676}{rad.s^{-1}} \times \qty{22.5}{m} = \qty{37.70}{m.s^{-1}}$.}
  \answer{(a) \final{$\omega = \qty{1.68}{rad.s^{-1}}$} \quad
    (b) \final{$v = \qty{271}{km.h^{-1}}$} \quad (c) the same $\omega$, but half the speed
    (\qty{37.7}{m.s^{-1}})}
  \checkstep{one turn takes $60 \div 16.0 = \qty{3.75}{s}$, and
    $\dfrac{2\pi r}{T} = \dfrac{2\pi \times 45.0}{3.75} = \qty{75.4}{m.s^{-1}}$ \cmark.}
  \begin{markingcriteria}
    \criterion{(a) Correct $\omega$, \qty{1.68}{rad.s^{-1}}}{1}
    \criterion{(b) Correct speed, \qty{271}{km.h^{-1}}}{1}
    \criterion{(c) Same $\omega$ and half the speed, with a reason}{1}
  \end{markingcriteria}
\end{practiceans}

% ---- Q12 satellite altitude from period
\begin{practiceans}
  \given{$T = 100.0 \times \qty{60}{s} = \qty{6000}{s}$; \
    $GM_\mathrm{E} = \num{6.67e-11} \times \num{6.0e24} = \qty{4.002e14}{m^3.s^{-2}}$}
  \find{the altitude $h$}
  \formula{$\dfrac{r^3}{T^2} = \dfrac{GM}{4\pi^2}$, so $r = \sqrt[3]{\dfrac{GMT^2}{4\pi^2}}$;
    \quad then $h = r - r_\mathrm{E}$}
  \substitute{\eqlines{r^3 &= \frac{\qty{4.002e14}{m^3.s^{-2}} \times (\qty{6000}{s})^2}{4\pi^2}\\
      &= \frac{\qty{1.441e22}{m^3}}{39.48} = \qty{3.649e20}{m^3}\\
      r &= \sqrt[3]{\qty{3.649e20}{m^3}} = \qty{7.146e6}{m} \quad\text{(cube root)}\\
      h &= \qty{7.146e6}{m} - \qty{6.371e6}{m} = \qty{7.75e5}{m}}}
  \answer{\final{$h = \qty{775}{km}$} ($\qty{7.75e5}{m}$)}
  \checkstep{low Earth orbits are \qtyrange{200}{2000}{km} up with periods of about
    \qtyrange{90}{130}{min} \cmark.}
  \begin{markingcriteria}
    \criterion{Changes $T$ to seconds and substitutes into $\dfrac{r^3}{T^2} = \dfrac{GM}{4\pi^2}$}{1}
    \criterion{Correct radius, \qty{7.15e6}{m}}{1}
    \criterion{Subtracts $r_\mathrm{E}$: altitude \qty{775}{km}}{1}
  \end{markingcriteria}
\end{practiceans}

% ---- Q13 forearm
\begin{practiceans}
  \given{dumbbell: \qty{6.0}{kg} at $r_\perp = \qty{0.34}{m}$; forearm: \qty{1.8}{kg} at
    $r_\perp = \qty{0.16}{m}$; biceps: $r_\perp = \qty{0.040}{m}$. All forces are vertical and
    the forearm is horizontal, so each force is at \ang{90} to the forearm.}
  \find{the biceps force $F_\mathrm{b}$}
  \formula{$\tau = r_\perp F$ \quad and, for no change in rotation,
    $\sum\tau_\mathrm{clockwise} = \sum\tau_\mathrm{anticlockwise}$}
  \substitute{The weights turn the forearm clockwise about the elbow. The biceps turns it
    anticlockwise.\par\nobreak
    \eqlines{\text{weights: } & 6.0 \times 9.8 = \qty{58.8}{N}, \quad 1.8 \times 9.8 = \qty{17.64}{N}\\
      \tau_\mathrm{clockwise} &= (\qty{0.34}{m})(\qty{58.8}{N}) + (\qty{0.16}{m})(\qty{17.64}{N})\\
      &= \qty{19.99}{N.m} + \qty{2.822}{N.m} = \qty{22.81}{N.m}\\
      (\qty{0.040}{m})\,F_\mathrm{b} &= \qty{22.81}{N.m} \quad\text{(anticlockwise = clockwise)}\\
      F_\mathrm{b} &= \frac{\qty{22.81}{N.m}}{\qty{0.040}{m}} = \qty{570.4}{N} \quad\text{(divide by \qty{0.040}{m})}}}
  \answer{\final{$F_\mathrm{b} = \qty{5.7e2}{N}$} upwards}
  \checkstep{the muscle pulls about 7.5 times harder than the total weight (\qty{76}{N}),
    because it acts very close to the pivot. Muscles often work like this \cmark.}
  \begin{markingcriteria}
    \criterion{Correct torques of the two weights about the elbow}{1}
    \criterion{Uses balanced torques: $0.040\,F_\mathrm{b} = $ total clockwise torque}{1}
    \criterion{Correct answer, \qty{5.7e2}{N}}{1}
  \end{markingcriteria}
\end{practiceans}

% ---- Q14 golf
\begin{practiceans}
  \given{$u = 126 \div 3.6 = \qty{35.0}{m.s^{-1}}$; \ $\theta = \ang{24}$; \
    $a_y = \qty{-9.8}{m.s^{-2}}$; level ground}
  \find{(a) the maximum height $h_\mathrm{max}$ \quad (b) the range $R$}
  \formula{$u_x = u\cos\theta$, \ $u_y = u\sin\theta$; \ $v_y^2 = u_y^2 + 2a_y s_y$; \
    $s_y = u_y t + \frac{1}{2}a_y t^2$; \ $s_x = u_x t$}
  \substitute{\eqlines{u_x &= 35.0\cos\ang{24} = \qty{31.97}{m.s^{-1}}, \qquad
      u_y = 35.0\sin\ang{24} = \qty{14.24}{m.s^{-1}}}\par
    \textbf{(a)} At the top, $v_y = 0$:\par\nobreak
    \eqlines{0 &= (\qty{14.24}{m.s^{-1}})^2 + 2(\qty{-9.8}{m.s^{-2}})\,h_\mathrm{max}\\
      (\qty{19.6}{m.s^{-2}})\,h_\mathrm{max} &= \qty{202.7}{m^2.s^{-2}} \quad\text{(add $19.6h_\mathrm{max}$ to both sides)}\\
      h_\mathrm{max} &= \frac{\qty{202.7}{m^2.s^{-2}}}{\qty{19.6}{m.s^{-2}}} = \qty{10.34}{m} \quad\text{(divide both sides by 19.6)}}\par
    \textbf{(b)} The ball lands at the launch height, so $s_y = 0$:\par\nobreak
    \eqlines{0 &= 14.24\,t - 4.9\,t^2 = t\,(14.24 - 4.9\,t) \quad\text{(take out $t$)}\\
      t &= \frac{14.24}{4.9} = \qty{2.905}{s} \quad\text{($t = 0$ is the launch)}\\
      R &= u_x t = \qty{31.97}{m.s^{-1}} \times \qty{2.905}{s} = \qty{92.89}{m}}}
  \answer{(a) \final{$h_\mathrm{max} = \qty{10.3}{m}$} \quad (b) \final{$R = \qty{92.9}{m}$}}
  \checkstep{the range card gives
    $R = \dfrac{u^2\sin 2\theta}{g} = \dfrac{35.0^2 \times \sin\ang{48}}{9.8} = \qty{92.9}{m}$ \cmark.}
  \begin{markingcriteria}
    \criterion{(a) Changes \si{km.h^{-1}} to \si{m.s^{-1}} and finds $u_y$ (and $u_x$)}{1}
    \criterion{(a) Correct maximum height, \qty{10.3}{m}}{1}
    \criterion{(b) Correct time of flight, \qty{2.91}{s}}{1}
    \criterion{(b) Correct range, \qty{92.9}{m}}{1}
  \end{markingcriteria}
\end{practiceans}

% ---- Q15 Mercury
\begin{practiceans}
  \given{$M = \qty{3.30e23}{kg}$, \ $r = \qty{2.44e6}{m}$}
  \find{(a) $v_\mathrm{esc}$ \quad (b) $\dfrac{v_\mathrm{esc}}{v_\text{orbital}}$}
  \formula{$v_\mathrm{esc} = \sqrt{\dfrac{2GM}{r}}$ (not on the formula sheet); in orbit,
    $\dfrac{GMm}{r^2} = \dfrac{mv^2}{r}$ gives $v_\text{orbital} = \sqrt{\dfrac{GM}{r}}$}
  \substitute{\textbf{(a)}\par\nobreak
    \eqlines{\frac{2GM}{r} &= \frac{2 \times \num{6.67e-11} \times \num{3.30e23}}{\num{2.44e6}}
        = \qty{1.804e7}{m^2.s^{-2}}\\
      v_\mathrm{esc} &= \sqrt{\qty{1.804e7}{m^2.s^{-2}}} = \qty{4.248e3}{m.s^{-1}}}\par
    \textbf{(b)} Divide the two formulas:\par\nobreak
    \eqlines{\frac{v_\mathrm{esc}}{v_\text{orbital}}
        &= \frac{\sqrt{2GM/r}}{\sqrt{GM/r}} = \sqrt{\frac{2GM/r}{GM/r}} = \sqrt{2} = 1.41
        \quad\text{($GM/r$ cancels)}}\par
    With numbers: $v_\text{orbital} = \sqrt{\qty{9.021e6}{m^2.s^{-2}}} = \qty{3.003e3}{m.s^{-1}}$,
    and $4248 \div 3003 = 1.41$.}
  \answer{(a) \final{$v_\mathrm{esc} = \qty{4.25e3}{m.s^{-1}}$} \quad (b) the speed must
    increase by a factor of $\sqrt{2} = 1.41$ \cmark}
  \checkstep{Mercury's escape velocity is much smaller than Earth's (\qty{1.12e4}{m.s^{-1}}),
    because Mercury has a much smaller mass \cmark.}
  \begin{markingcriteria}
    \criterion{(a) Correct substitution into $v_\mathrm{esc} = \sqrt{\dfrac{2GM}{r}}$}{1}
    \criterion{(a) Correct answer, \qty{4.25e3}{m.s^{-1}}}{1}
    \criterion{(b) Shows that $v_\mathrm{esc} \div v_\text{orbital} = \sqrt{2}$ (by algebra or
      with numbers)}{1}
  \end{markingcriteria}
\end{practiceans}

% ---- Q16 swing ride
\begin{practiceans}
  \given{$d = \qty{4.0}{m}$, \ $L = \qty{5.0}{m}$, \ $\theta = \ang{40}$ from the vertical, \
    $g = \qty{9.8}{m.s^{-2}}$}
  \find{(a) $r$ \quad (b) $v$ \quad (c) the period $T$}
  \formula{$r = d + L\sin\theta$; \ $\tan\theta = \dfrac{v^2}{rg}$ (not on the formula sheet);
    \ $v = \dfrac{2\pi r}{T}$}
  \substitute{\textbf{(a)} The circle's centre is on the axis, so $r$ is the arm plus the
    sideways part of the chain:
    $r = \qty{4.0}{m} + \qty{5.0}{m} \times \sin\ang{40} = \qty{4.0}{m} + \qty{3.214}{m}
    = \qty{7.214}{m}$.\par
    \textbf{(b)} Two forces act on the rider: the chain tension $F_\mathrm{T}$ (along the chain)
    and the weight $mg$.\par\nobreak
    \eqlines{F_\mathrm{T}\cos\theta &= mg \quad\text{(vertical: forces balance)}\\
      F_\mathrm{T}\sin\theta &= \frac{mv^2}{r} \quad\text{(horizontal: net force to the centre)}\\
      \tan\theta &= \frac{v^2}{rg} \quad\text{(divide the second line by the first)}\\
      v^2 &= rg\tan\theta = \qty{7.214}{m} \times \qty{9.8}{m.s^{-2}} \times 0.8391
        = \qty{59.32}{m^2.s^{-2}}\\
      v &= \sqrt{\qty{59.32}{m^2.s^{-2}}} = \qty{7.702}{m.s^{-1}}}\par
    \textbf{(c)} $T = \dfrac{2\pi r}{v} = \dfrac{2\pi \times \qty{7.214}{m}}{\qty{7.702}{m.s^{-1}}}
    = \dfrac{\qty{45.33}{m}}{\qty{7.702}{m.s^{-1}}} = \qty{5.885}{s}$}
  \answer{(a) \final{$r = \qty{7.2}{m}$} \quad (b) \final{$v = \qty{7.7}{m.s^{-1}}$} \quad
    (c) \final{$T = \qty{5.9}{s}$}}
  \checkstep{units: $\sqrt{\unit{m} \times \unit{m.s^{-2}}} = \unit{m.s^{-1}}$ \cmark. If you
    forget the arm $d$, you get $r = \qty{3.2}{m}$ and $v = \qty{5.1}{m.s^{-1}}$, which is too
    small.}
  \begin{markingcriteria}
    \criterion{(a) Correct radius, \qty{7.2}{m} (includes the arm length $d$)}{1}
    \criterion{(b) Uses the force components to get $\tan\theta = \dfrac{v^2}{rg}$ (or
      equivalent)}{1}
    \criterion{(b) Correct speed, \qty{7.7}{m.s^{-1}}}{1}
    \criterion{(c) Correct period, \qty{5.9}{s}}{1}
  \end{markingcriteria}
\end{practiceans}

% ---- Q17 vertical circle
\begin{practiceans}
  \given{$m = \qty{0.20}{kg}$, \ $r = \qty{0.80}{m}$, \ $v_\mathrm{top} = \qty{3.5}{m.s^{-1}}$,
    \ $mg = 0.20 \times 9.8 = \qty{1.96}{N}$}
  \find{(a) $F_\mathrm{T}$ at the top \quad (b) $v_\mathrm{bottom}$ \quad (c) $F_\mathrm{T}$ at
    the bottom}
  \formula{top: $F_\mathrm{T} + mg = \dfrac{mv^2}{r}$; \ bottom: $F_\mathrm{T} - mg = \dfrac{mv^2}{r}$
    (not on the formula sheet); \
    energy: $\tfrac{1}{2}mv_\mathrm{bottom}^2 = \tfrac{1}{2}mv_\mathrm{top}^2 + mg(2r)$}
  \substitute{\textbf{(a)} At the top, the tension and the weight both point down, towards
    the centre.\par\nobreak
    \eqlines{\frac{mv^2}{r} &= \frac{\qty{0.20}{kg} \times (\qty{3.5}{m.s^{-1}})^2}{\qty{0.80}{m}}
        = \qty{3.063}{N}\\
      F_\mathrm{T} &= \qty{3.063}{N} - \qty{1.96}{N} = \qty{1.103}{N}
        \quad\text{(subtract $mg$ from both sides)}}\par
    \textbf{(b)} The bottom is $2r = \qty{1.60}{m}$ below the top. The ball loses
    gravitational potential energy $mg(2r)$ and gains the same amount of kinetic energy.
    Cancel $m$ and multiply by 2:\par\nobreak
    \eqlines{v_\mathrm{bottom}^2 &= v_\mathrm{top}^2 + 4gr
        = \qty{12.25}{m^2.s^{-2}} + 4 \times \qty{9.8}{m.s^{-2}} \times \qty{0.80}{m}\\
      &= 12.25 + 31.36 = \qty{43.61}{m^2.s^{-2}}\\
      v_\mathrm{bottom} &= \sqrt{\qty{43.61}{m^2.s^{-2}}} = \qty{6.604}{m.s^{-1}}}\par
    \textbf{(c)} At the bottom, the tension points up (to the centre) and the weight points
    down.\par\nobreak
    \eqlines{\frac{mv^2}{r} &= \frac{\qty{0.20}{kg} \times \qty{43.61}{m^2.s^{-2}}}{\qty{0.80}{m}}
        = \qty{10.90}{N}\\
      F_\mathrm{T} &= \qty{10.90}{N} + \qty{1.96}{N} = \qty{12.86}{N}
        \quad\text{(add $mg$ to both sides)}}}
  \answer{(a) \final{$F_\mathrm{T} = \qty{1.1}{N}$} \quad
    (b) \final{$v = \qty{6.6}{m.s^{-1}}$} \quad (c) \final{$F_\mathrm{T} = \qty{12.9}{N}$}}
  \checkstep{the top speed is above the smallest possible top speed,
    $\sqrt{gr} = \qty{2.8}{m.s^{-1}}$, so the string stays tight ($F_\mathrm{T} > 0$) \cmark. At
    the bottom the string must hold up the ball \emph{and} pull it to the centre, so the
    tension is much bigger \cmark.}
  \begin{markingcriteria}
    \criterion{(a) Writes $F_\mathrm{T} + mg = \dfrac{mv^2}{r}$ (both forces towards the centre)}{1}
    \criterion{(a) Correct tension at the top, \qty{1.1}{N}}{1}
    \criterion{(b) Correct speed at the bottom, \qty{6.6}{m.s^{-1}}, using energy (height change $2r$)}{1}
    \criterion{(c) Correct tension at the bottom, \qty{12.9}{N}}{1}
  \end{markingcriteria}
\end{practiceans}

% ---- Q18 orbit raise
\begin{practiceans}
  \given{$m = \qty{900}{kg}$; \ $r_1 = \qty{6.371e6}{m} + \qty{0.350e6}{m} = \qty{6.721e6}{m}$;
    \ $T_2 = 4.00 \times \qty{3600}{s} = \qty{14400}{s}$; \ $GM_\mathrm{E} = \qty{4.002e14}{m^3.s^{-2}}$}
  \find{(a) $r_2$ \quad (b) the energy needed, $\Delta E$ \quad (c) $\Delta K$}
  \formula{$r = \sqrt[3]{\dfrac{GMT^2}{4\pi^2}}$; \
    $\Delta E = \dfrac{GMm}{2}\left(\dfrac{1}{r_1} - \dfrac{1}{r_2}\right)$ (not on the formula
    sheet); \ $K = \dfrac{GMm}{2r}$, so $\Delta K = -\Delta E$}
  \substitute{\textbf{(a)}\par\nobreak
    \eqlines{r_2^3 &= \frac{\qty{4.002e14}{m^3.s^{-2}} \times (\qty{14400}{s})^2}{4\pi^2}\\
      &= \frac{\qty{8.299e22}{m^3}}{39.48} = \qty{2.102e21}{m^3}\\
      r_2 &= \sqrt[3]{\qty{2.102e21}{m^3}} = \qty{1.281e7}{m}}\par
    \textbf{(b)}\par\nobreak
    \eqlines{GMm &= \qty{4.002e14}{m^3.s^{-2}} \times \qty{900}{kg} = \qty{3.602e17}{J.m}\\
      \frac{1}{r_1} - \frac{1}{r_2} &= \frac{1}{\qty{6.721e6}{m}} - \frac{1}{\qty{1.281e7}{m}}\\
      &= \qty{1.488e-7}{m^{-1}} - \qty{7.806e-8}{m^{-1}} = \qty{7.072e-8}{m^{-1}}\\
      \Delta E &= \frac{\qty{3.602e17}{J.m}}{2} \times \qty{7.072e-8}{m^{-1}} = \qty{1.274e10}{J}}\par
    \textbf{(c)} $\Delta K = \dfrac{GMm}{2r_2} - \dfrac{GMm}{2r_1}
    = -\dfrac{GMm}{2}\left(\dfrac{1}{r_1} - \dfrac{1}{r_2}\right) = -\Delta E
    = \qty{-1.274e10}{J}$}
  \answer{(a) \final{$r_2 = \qty{1.28e7}{m}$} \quad (b) \final{$\Delta E = \qty{1.27e10}{J}$}\par
    (c) \final{$\Delta K = \qty{-1.27e10}{J}$}: the kinetic energy decreases, so the
    satellite moves \textbf{more slowly} in the new orbit (\qty{7.72}{km.s^{-1}} becomes
    \qty{5.59}{km.s^{-1}}).}
  \checkstep{$E_1 = -\dfrac{GMm}{2r_1} = \qty{-2.680e10}{J}$ and
    $E_2 = \qty{-1.406e10}{J}$; $E_2 - E_1 = \qty{1.27e10}{J}$ \cmark. $U$ rises by twice this,
    \qty{2.55e10}{J}, so $\Delta K + \Delta U = \Delta E$ \cmark.}
  \begin{markingcriteria}
    \criterion{(a) Correct radius of the new orbit, \qty{1.28e7}{m}}{1}
    \criterion{(b) Uses $r_1 = r_\mathrm{E} + \text{altitude}$ in the total-energy change}{1}
    \criterion{(b) Correct energy, \qty{1.27e10}{J}}{1}
    \criterion{(c) $\Delta K = \qty{-1.27e10}{J}$, so the satellite is slower}{1}
  \end{markingcriteria}
\end{practiceans}

% ===================================================================== Section III
% ---- Q19 work and energy, horizontal vs vertical circle
\begin{practiceans}
  \begin{modelanswer}
  In both situations the tension pulls along the string, towards the centre. The velocity is
  along the tangent. So the tension is at \ang{90} to the motion at every instant. Its work is
  $W = Fs\cos\ang{90} = 0$. The tension changes only the \emph{direction} of the velocity.

  \textbf{Situation A.} The weight and the normal force from the table are vertical. The ball
  moves horizontally, so these forces are also at \ang{90} to the motion and do no work. No
  force does work, so $W_\mathrm{net} = 0$. By the work--energy theorem,
  $\Delta K = W_\mathrm{net} = 0$. The kinetic energy stays the same, so the speed stays
  constant.

  \textbf{Situation B.} The ball moves up and down, and its weight is vertical. Except at the
  very top and bottom, the weight has a component along the path, so gravity does work. On the
  way up, gravity does negative work: kinetic energy changes into potential energy, and the
  ball slows down. On the way down, gravity does positive work: potential energy changes into
  kinetic energy, and the ball speeds up. The total energy $K + U$ stays constant (no friction
  or air resistance). So the speed is smallest at the top and largest at the bottom.
  \end{modelanswer}
  \begin{markingcriteria}
    \criterion{Explains that the tension is at \ang{90} to the velocity, so it does no work in
      either situation; explains that in A no force does work, so $K$ and the speed are
      constant; explains that in B gravity has a component along the path and does work,
      changing $K$ into $U$ and back while $K + U$ stays constant}{4}
    \criterion{Explains why the speed is constant in A and why it changes in B, with one
      point missing or unclear}{3}
    \criterion{Explains one situation correctly, OR states that the tension does no work and
      that gravity does work in B}{2}
    \criterion{States one relevant point, for example ``a force at \ang{90} to the motion does
      no work''}{1}
  \end{markingcriteria}
\end{practiceans}

% ---- Q20 LEO vs GEO
\begin{practiceans}
  \begin{modelanswer}
  \textbf{Similarities.} In both orbits, gravity provides the centripetal force. Both obey
  Kepler's third law, $\dfrac{r^3}{T^2} = \dfrac{GM}{4\pi^2}$, so the higher orbit has the
  longer period and the lower speed.

  \textbf{Differences.}
  \begin{itemize}
    \item \textbf{Altitude:} a low Earth orbit (LEO) is about \qtyrange{200}{2000}{km} above
          the surface. A geostationary orbit (GEO) is about \qty{36000}{km} above the equator.
    \item \textbf{Period:} a LEO satellite orbits in about \qtyrange{90}{130}{min}. A GEO
          satellite takes one day (about \qty{24}{h}), the same time as the Earth takes to spin. It
          moves west to east, above the equator, so it stays above one point on the ground.
    \item \textbf{Speed:} about \qtyrange{6.9}{7.8}{km.s^{-1}} in LEO, but only \qty{3.1}{km.s^{-1}} in
          GEO.
  \end{itemize}

  \textbf{Uses.} LEO satellites are close to the Earth, so their cameras see fine detail (for
  example mapping and imaging), and signals take only milliseconds. But each one crosses the
  sky in minutes, so many are needed for constant cover, and air drag slowly lowers them. A
  GEO satellite stays above one point, so a dish on the ground can point at it without moving.
  This suits TV and communications, and weather satellites that watch the same half of the
  Earth all the time. But it is far away, so its images show less detail and its signals are
  delayed by about \qty{0.24}{s}.
  \end{modelanswer}
  \begin{markingcriteria}
    \criterion{Gives a similarity and differences in altitude AND period, with approximate
      values; links one use of each orbit to a property of that orbit}{4}
    \criterion{Gives differences in altitude and period, with values; links one use to a
      property}{3}
    \criterion{Gives a difference in altitude or period AND one correct use; OR gives
      differences but no uses}{2}
    \criterion{States one relevant feature of either orbit}{1}
  \end{markingcriteria}
\end{practiceans}

% ---- Q21 evaluate the stopper experiment
\begin{practiceans}
  \begin{modelanswer}
  \textbf{Judgement.} As written, the method has low validity and low reliability. With the
  changes below, it would be a valid and reliable test of $F_\mathrm{c} \propto v^2$.

  \textbf{Validity} (does it fairly test $F_\mathrm{c} \propto v^2$?)
  \begin{itemize}[beginpenalty=10000]
    \item The radius is not controlled. Nothing marks the string, so the radius can change
          from trial to trial. It is measured only once, at the end, so $v = \dfrac{2\pi r}{T}$
          uses a wrong $r$ for most trials. \emph{Improvement:} clip a paper clip to the string
          just below the tube and keep it there while whirling. Measure $r$ from the top of
          the tube to the centre of the stopper before each trial.
    \item The force on the stopper may not be $Mg$. The tension equals $Mg$ only while the
          hanging mass stays still. If the mass moves up or down, or if the string rubs on
          the tube (friction), the pull on the stopper is not $Mg$. Whirling ``above your
          head'' also gives no check that the circle is horizontal. \emph{Improvement:} use a
          smooth tube, and whirl the stopper in a horizontal circle at a speed that keeps the
          clip (and so the hanging mass) still.
  \end{itemize}
  (A small sag of the string below the horizontal is \emph{not} a problem. If you use the
  string length for $r$ in $v = \dfrac{2\pi r}{T}$, the sag cancels out. See the
  whirling-stopper experiment in \cref{ch2:sec:exp}.)

  \textbf{Reliability} (would repeats give the same results?)
  \begin{itemize}[beginpenalty=10000]
    \item One revolution takes less than a second. A reaction time of about \qty{0.2}{s} is
          a large part of this, so each value of $T$ has a large random error. There is only
          one trial for each mass. \emph{Improvement:} time 20 revolutions and divide by 20.
          Repeat each mass three times and use the mean. (A video or a light gate can also
          time the turns.)
  \end{itemize}

  \textbf{Analysis.} Calculate $v$ and $v^2$ for each mass. Plot $F_\mathrm{c} = Mg$ against
  $v^2$. A straight line through the origin supports $F_\mathrm{c} \propto v^2$.
  \end{modelanswer}
  \begin{markingcriteria}
    \criterion{Makes a clear judgement about validity AND reliability; identifies at least
      three specific problems, explains how each affects the results, and gives a matching
      improvement for each}{5}
    \criterion{Identifies three problems with improvements and refers to validity and
      reliability; the judgement or one explanation is weak}{4}
    \criterion{Identifies two problems, explains their effect and gives improvements}{3}
    \criterion{Identifies two problems, OR one problem with its effect and an improvement}{2}
    \criterion{Identifies one relevant problem or improvement}{1}
  \end{markingcriteria}
\end{practiceans}

% ---- Q22 spring launcher graph
\begin{practiceans}
  \textbf{(a)} $\sin 2\theta$ (calculator in degree mode; for \ang{10}, $2\theta = \ang{20}$):
  \begin{center}
  \begin{tabular}{@{}lcccccc@{}}
    \toprule
    $\theta$ & \ang{10} & \ang{15} & \ang{20} & \ang{25} & \ang{30} & \ang{45}\\
    $\sin 2\theta$ & 0.342 & 0.500 & 0.643 & 0.766 & 0.866 & 1.000\\
    $R$ (\si{m}) & 0.660 & 0.958 & 1.236 & 1.482 & 1.661 & 1.918\\
    \bottomrule
  \end{tabular}
  \end{center}

  \textbf{(b)} See the graph below. The points lie close to a straight line through the
  origin, as the model predicts.
  \begin{figure}
    \centering
    \begin{tikzpicture}
      \begin{axis}[hsc, width=0.78\linewidth, height=0.44\linewidth,
          xmin=0, xmax=1.1, ymin=0, ymax=2.2,
          xlabel={$\sin 2\theta$}, ylabel={$R$ (\si{m})},
          xtick={0,0.2,...,1.0}, ytick={0,0.4,...,2.0},
          xticklabel style={/pgf/number format/.cd, fixed, fixed zerofill, precision=1},
          yticklabel style={/pgf/number format/.cd, fixed, fixed zerofill, precision=1}]
        \addplot[hscseriesB, line width=1pt, domain=0:1.08, samples=2, mark=none] {1.92*x};
        \addplot[hscseriesA, only marks, mark=*, mark size=2pt]
          coordinates {(0.342,0.660) (0.500,0.958) (0.643,1.236) (0.766,1.482)
                       (0.866,1.661) (1.000,1.918)};
        \draw[hscmuted, densely dotted, line width=0.7pt] (axis cs:0,0) -- (axis cs:1.0,0)
          -- (axis cs:1.0,1.92);
        \node[anchor=south, font=\footnotesize, text=hscmuted] at (axis cs:0.80,0.02)
          {run $= 1.000$};
        \node[anchor=west, font=\footnotesize, text=hscmuted] at (axis cs:1.01,0.75)
          {rise $= \qty{1.92}{m}$};
        \node[anchor=south east, font=\footnotesize, text=hscseriesB] at (axis cs:0.45,1.15)
          {line of best fit};
      \end{axis}
    \end{tikzpicture}
    \caption*{Answer to \cref{pq:launcher}(b): $R$ against $\sin 2\theta$.}
  \end{figure}

  \textbf{(c)}
  \given{two points on the line of best fit: $(0, 0)$ and $(1.000, \qty{1.92}{m})$;
    $g = \qty{9.8}{m.s^{-2}}$}
  \find{the gradient, then the launch speed $u$}
  \formula{$R = \left(\dfrac{u^2}{g}\right)\sin 2\theta$ has the form $y = kx$, so
    gradient $= \dfrac{u^2}{g}$}
  \substitute{\eqlines{\text{gradient} &= \frac{\text{rise}}{\text{run}}
        = \frac{\qty{1.92}{m} - 0}{1.000 - 0} = \qty{1.92}{m}\\
      u^2 &= g \times \text{gradient} \quad\text{(multiply both sides by $g$)}\\
      &= \qty{9.8}{m.s^{-2}} \times \qty{1.92}{m} = \qty{18.82}{m^2.s^{-2}}\\
      u &= \sqrt{\qty{18.82}{m^2.s^{-2}}} = \qty{4.338}{m.s^{-1}}}}
  \answer{\final{$u = \qty{4.34}{m.s^{-1}}$} (any value from about \qtyrange{4.30}{4.37}{m.s^{-1}},
    from a good line, is fine)}
  \checkstep{$\sin 2\theta$ has no unit, so the gradient has the unit of $R$ (\si{m}), and
    $\sqrt{\unit{m.s^{-2}} \times \unit{m}} = \unit{m.s^{-1}}$ \cmark.}

  \textbf{(d)} With $u = \qty{4.40}{m.s^{-1}}$, every range would be 2 to 3\,\% longer
  than the measured range. So something makes every range a little shorter. Any one
  sensible reason is enough:
  \begin{itemize}
    \item The real launch speed is lower than the maker says. For example, the spring is
          weaker than when it was new, or friction inside the launcher takes some energy.
    \item Air resistance on the real ball slows it down and shortens every range.
  \end{itemize}
  Both are systematic errors: they make the gradient, and the speed found from it, a little
  too small.
  \begin{markingcriteria}
    \criterion{(a) All six values of $\sin 2\theta$ correct}{1}
    \criterion{(b) All points plotted correctly}{1}
    \criterion{(b) Straight line of best fit, through (or very near) the origin}{1}
    \criterion{(c) Gradient from the line, about \qty{1.92}{m}, with gradient $= \dfrac{u^2}{g}$}{1}
    \criterion{(c) Correct launch speed, about \qty{4.34}{m.s^{-1}}}{1}
    \criterion{(d) A sensible reason, for example the real launch speed is lower (weaker
      spring or friction in the launcher), or air resistance shortens every range}{1}
  \end{markingcriteria}
\end{practiceans}

% ---- Q23 TRAPPIST-1
\begin{practiceans}
  \textbf{(a)} The star's gravity provides the centripetal force on the planet ($m$):
  \begin{align*}
    \frac{GMm}{r^2} &= \frac{mv^2}{r} && \text{gravity = centripetal force}\\
    \frac{GM}{r} &= v^2 && \text{cancel $m$; multiply both sides by $r$}\\
    \frac{GM}{r} &= \frac{4\pi^2 r^2}{T^2} && \text{put in $v = \dfrac{2\pi r}{T}$, so $v^2 = \dfrac{4\pi^2r^2}{T^2}$}\\
    \frac{r^3}{T^2} &= \frac{GM}{4\pi^2} && \text{multiply both sides by $\dfrac{r}{4\pi^2}$}
  \end{align*}

  \textbf{(b)}
  \given{two points on the line: $(0, 0)$ and $(T^2, r^3) = (\qty{2.50e12}{s^2}, \qty{7.56e29}{m^3})$;
    $G = \qty{6.67e-11}{N.m^2.kg^{-2}}$}
  \find{the gradient, then the star's mass $M$}
  \formula{gradient $= \dfrac{r^3}{T^2} = \dfrac{GM}{4\pi^2}$, \ so \
    $M = \dfrac{4\pi^2 \times \text{gradient}}{G}$}
  \substitute{\eqlines{\text{gradient} &= \frac{\qty{7.56e29}{m^3} - 0}{\qty{2.50e12}{s^2} - 0}
        = \qty{3.024e17}{m^3.s^{-2}}\\
      M &= \frac{4\pi^2 \times \qty{3.024e17}{m^3.s^{-2}}}{\qty{6.67e-11}{N.m^2.kg^{-2}}}\\
      &= \frac{\qty{1.194e19}{m^3.s^{-2}}}{\qty{6.67e-11}{N.m^2.kg^{-2}}}
        = \qty{1.790e29}{kg}}}
  \answer{gradient \final{\qty{3.0e17}{m^3.s^{-2}}}; \ \final{$M = \qty{1.8e29}{kg}$}}
  \checkstep{this is about 9\,\% of the Sun's mass (\qty{2.0e30}{kg}). TRAPPIST-1 is a small,
    cool star \cmark. Read the powers of ten from the axis labels: $7.56$ means
    $\qty{7.56e29}{m^3}$.}

  \textbf{(c)} All the planets orbit the same star, so they share the same gradient:\par\nobreak
  \eqlines{T^2 &= \frac{r^3}{\text{gradient}} = \frac{(\qty{7.01e9}{m})^3}{\qty{3.024e17}{m^3.s^{-2}}}
      = \frac{\qty{3.445e29}{m^3}}{\qty{3.024e17}{m^3.s^{-2}}} = \qty{1.139e12}{s^2}\\
    T &= \sqrt{\qty{1.139e12}{s^2}} = \qty{1.067e6}{s} \quad\text{(square root)}\\
    T &= \frac{\qty{1.067e6}{s}}{\qty{86400}{s}} = 12.4\ \text{days} \quad\text{(one day = \qty{86400}{s})}}\par
  Answer: \final{$T = 12.4$ days}. (The measured period of planet g is 12.35 days.)
  \begin{markingcriteria}
    \criterion{(a) States that gravity provides the centripetal force:
      $\dfrac{GMm}{r^2} = \dfrac{mv^2}{r}$}{1}
    \criterion{(a) Substitutes $v = \dfrac{2\pi r}{T}$ and rearranges to
      $\dfrac{r^3}{T^2} = \dfrac{GM}{4\pi^2}$}{1}
    \criterion{(b) Gradient from two points on the line, about \qty{3.0e17}{m^3.s^{-2}}}{1}
    \criterion{(b) Uses $M = \dfrac{4\pi^2 \times \text{gradient}}{G}$}{1}
    \criterion{(b) Correct mass, about \qty{1.8e29}{kg}}{1}
    \criterion{(c) Correct period, about 12.4 days}{1}
  \end{markingcriteria}
\end{practiceans}

% ---- Q24 elliptical orbit
\begin{practiceans}
  \textbf{(a)}
  \given{$r_\mathrm{p} = \qty{6.90e6}{m}$, \ $v_\mathrm{p} = \qty{9.63e3}{m.s^{-1}}$, \
    $r_\mathrm{a} = \qty{2.76e7}{m}$}
  \find{the speed at apogee, $v_\mathrm{a}$}
  \formula{Kepler's second law at the closest and farthest points:
    $r_\mathrm{p}v_\mathrm{p} = r_\mathrm{a}v_\mathrm{a}$ (not on the formula sheet)}
  \substitute{\eqlines{v_\mathrm{a} &= \frac{r_\mathrm{p}v_\mathrm{p}}{r_\mathrm{a}}
      \quad\text{(divide both sides by $r_\mathrm{a}$)}\\
      &= \frac{\qty{6.90e6}{m} \times \qty{9.63e3}{m.s^{-1}}}{\qty{2.76e7}{m}}
      = \qty{2.408e3}{m.s^{-1}}}}
  \answer{\final{$v_\mathrm{a} = \qty{2.41e3}{m.s^{-1}}$}}
  \checkstep{$r_\mathrm{a}$ is 4 times $r_\mathrm{p}$, so the speed is 4 times smaller \cmark.}

  \textbf{(b)}
  \given{$m = \qty{400}{kg}$; \
    $GMm = \qty{4.002e14}{m^3.s^{-2}} \times \qty{400}{kg} = \qty{1.601e17}{J.m}$}
  \find{$E = K + U$ at perigee (subscript p) and at apogee (subscript a)}
  \formula{$E = K + U = \dfrac{1}{2}mv^2 - \dfrac{GMm}{r}$}
  \substitute{\eqlines{K_\mathrm{p} &= \tfrac{1}{2} \times \qty{400}{kg} \times (\qty{9.63e3}{m.s^{-1}})^2
        = \qty{1.855e10}{J}\\
      U_\mathrm{p} &= -\frac{\qty{1.601e17}{J.m}}{\qty{6.90e6}{m}} = \qty{-2.320e10}{J}\\
      E_\mathrm{p} &= \qty{1.855e10}{J} + (\qty{-2.320e10}{J}) = \qty{-4.65e9}{J}}\par
    \eqlines{K_\mathrm{a} &= \tfrac{1}{2} \times \qty{400}{kg} \times (\qty{2.408e3}{m.s^{-1}})^2
        = \qty{1.160e9}{J}\\
      U_\mathrm{a} &= -\frac{\qty{1.601e17}{J.m}}{\qty{2.76e7}{m}} = \qty{-5.800e9}{J}\\
      E_\mathrm{a} &= \qty{1.160e9}{J} + (\qty{-5.800e9}{J}) = \qty{-4.64e9}{J}}}
  \answer{\final{$E = \qty{-4.6e9}{J}$} at both points.}
  \textbf{Comment:} the total energy is the same at perigee and apogee. (The tiny
  difference comes from rounding.) Only gravity acts on the satellite, so its mechanical
  energy is conserved: kinetic energy changes into potential energy and back. $E$ is
  negative, so the satellite is bound to the Earth: it cannot escape.

  \textbf{(c)}
  \begin{modelanswer}
  \textbf{Forces.} Gravity always pulls the satellite towards the Earth's centre. Between
  perigee and apogee, the satellite moves away from the Earth. So its velocity is not at
  \ang{90} to gravity: gravity has a component that points \emph{backwards}, against the
  velocity. This component slows the satellite down. Gravity does negative work on it.

  \textbf{Energy.} As $r$ increases, $U = -\dfrac{GMm}{r}$ increases (it becomes less
  negative). The total energy $E = K + U$ stays constant, so $K$ must decrease by the same
  amount: here $U$ rises by \qty{1.74e10}{J} and $K$ falls by \qty{1.74e10}{J}. A smaller $K$
  means a smaller speed. This agrees with Kepler's second law: the satellite moves fastest at
  perigee and slowest at apogee.
  \end{modelanswer}
  \begin{markingcriteria}
    \criterion{(a) Correct speed at apogee, \qty{2.41e3}{m.s^{-1}}}{1}
    \criterion{(b) Correct total energy at perigee, \qty{-4.65e9}{J}}{1}
    \criterion{(b) Correct total energy at apogee, \qty{-4.64e9}{J}}{1}
    \criterion{(b) States that $E$ is the same (conserved), because only gravity acts}{1}
    \criterion{(c) Explains with forces (a component of gravity opposes the velocity, so
      gravity does negative work) AND with energy ($U$ increases and $E$ is constant, so $K$
      and the speed decrease)}{3}
    \criterion{(c) Explains correctly with forces OR with energy}{2}
    \criterion{(c) States one relevant point}{1}
  \end{markingcriteria}
\end{practiceans}

\end{practiceanswers}
